% ============================================================
%  chapters/ch27-yarncrawler-systems.tex
% ============================================================

\chapter{Yarncrawler Systems}

\chapterepigraph{%
  A city that repairs itself\\
  is not merely convenient.\\
  It is epistemically honest about its own entropy.
}{Flyxion, \textit{Semantic Infrastructure}}

\sidenote{Yarncrawler:\\adaptive repair\\network;\\ connects to\\App.\,O on\\civilization\\dynamics}

Chapter~26 developed constraint ecology as the governing dynamics of
living knowledge systems. This chapter extends the framework to the
engineering question: how do we build systems that \emph{maintain}
their constraint ecology over time — that repair entropy, adapt to
change, and preserve accessibility without requiring centralized
intervention?

The Yarncrawler framework is the answer. It is a model of adaptive
maintenance: distributed systems that move through their environment
repairing, connecting, and reorganizing rather than building from scratch.

\section{Repair Networks}

Physical infrastructure degrades. Roads develop potholes; pipes corrode;
cables fray. The conventional response is scheduled maintenance:
periodic inspection and repair by specialized crews dispatched from
central depots. This is a centralized, reactive, low-frequency approach.

Yarncrawler systems propose an alternative: continuous, distributed,
proactive maintenance by small autonomous agents that are always
embedded in the infrastructure they maintain.

\begin{definition}{Repair Network}{repair-network}
  A \emph{repair network} is a tuple $(I, R, \tau_R, \kappa_I)$ where:
  \begin{itemize}
    \item $I$ is the infrastructure being maintained (a graph of
          nodes and edges),
    \item $R = \{r_1, \ldots, r_m\}$ is a collection of repair
          agents (Yarncrawlers),
    \item $\tau_R : I \times R \to I$ is the repair function
          (how agent $r$ at location $\ell$ modifies $I$),
    \item $\kappa_I : I \to \mathbb{R}_{\geq 0}$ is the
          infrastructure constraint violation functional (how
          degraded the infrastructure is at each point).
  \end{itemize}
\end{definition}

The goal of the repair network is to minimize $\int_I \kappa_I\,d\mu$
(total constraint violation) while each agent operates only with local
information — it can only observe the state of its immediate neighborhood.

\section{Moving Infrastructure}

A key insight of the Yarncrawler framework is that infrastructure need
not be static. The conventional picture separates the infrastructure
(static: roads, buildings, pipes) from the agents that use it (dynamic:
vehicles, people, water). The Yarncrawler picture blurs this distinction:
the repair agents \emph{are} a component of the infrastructure, and they
move.

\begin{definition}{Moving Infrastructure}{moving-infra}
  A \emph{moving infrastructure element} is an agent $r \in R$ that:
  \begin{enumerate}
    \item Provides infrastructure services (connectivity, resource
          transport, information relay) as it moves,
    \item Collects infrastructure state information (degradation level,
          resource levels, congestion) as it passes through,
    \item Modifies the infrastructure (repairs, reroutes, reinforces)
          in response to observed state.
  \end{enumerate}
\end{definition}

A Yarncrawler is a moving infrastructure element. As it traverses
the city, road, or network, it simultaneously serves the functions of
sensor, actor, and carrier. The repair is not a separate activity from
the service — it is woven into the service.

\section{Adaptive Maintenance}

Standard maintenance is scheduled: it occurs at fixed intervals regardless
of actual degradation state. Adaptive maintenance responds to the actual
state of the infrastructure.

\begin{definition}{Adaptive Maintenance Strategy}{adaptive-maint}
  A \emph{adaptive maintenance strategy} for Yarncrawler $r$ is a
  policy $\pi_r : \mathcal{O}_r \to \mathcal{A}_r$ mapping local
  observations to actions, where:
  \begin{itemize}
    \item $\mathcal{O}_r$ includes the current degradation level
          $\kappa_I(\ell_r)$ and neighboring degradation levels,
    \item $\mathcal{A}_r$ includes movement actions (which direction
          to move), repair actions (how much resource to expend on
          current location), and relay actions (what information to
          broadcast).
  \end{itemize}
  The strategy is \emph{accessibility-optimal} if $\pi_r$ maximizes
  the expected accessibility entropy of the infrastructure:
  \[
    \pi_r^* = \arg\max_{\pi_r} \mathbb{E}\bigl[S_I(t+H) \mid
    \mathcal{O}_r(t), \pi_r\bigr]
  \]
  over a horizon $H$.
\end{definition}

Accessibility-optimal maintenance prioritizes repairs that maximize
future infrastructure flexibility — not simply the most degraded
locations, but the locations where repair would most increase the
admissible uses of the infrastructure.

\section{Stigmergic Planning}

Yarncrawlers coordinate not through central planning but through
\emph{stigmergy}: indirect coordination through shared environment
modification.

\begin{definition}{Stigmergic Coordination}{stigmergy}
  Yarncrawlers use \emph{stigmergic coordination} when:
  \begin{enumerate}
    \item Each agent modifies the environment as it acts
          (by repairing, by leaving chemical/digital markers,
          by updating shared maps),
    \item Other agents' behavior is modified by these environmental
          changes without direct communication,
    \item The collective behavior that emerges is globally coherent
          despite only local information being available to any agent.
  \end{enumerate}
\end{definition}

Stigmergy is the coordination mechanism of ant colonies, termite
mounds, slime molds, and Wikipedia. Each agent acts on local information;
the environment records the aggregate effect; subsequent agents respond
to the recorded aggregate. The environment serves as the coordination
medium.

In the Yarncrawler framework, the infrastructure itself is the stigmergic
medium. Each repair leaves a trace (a physical trace — the repaired
surface — and potentially a digital trace in a maintenance log).
Other Yarncrawlers observe these traces and preferentially visit
recently repaired versus never-visited versus frequently-visited locations,
creating an emergent maintenance pattern that reflects the actual
degradation dynamics of the infrastructure.

\section{Connection to Semantic Infrastructure}

The Yarncrawler framework extends naturally from physical infrastructure
to semantic infrastructure (Chapter~25).

In semantic infrastructure, the ``degradation'' of a knowledge module
corresponds to:
\begin{itemize}
  \item Propositions becoming outdated (empirical facts changing),
  \item Connections breaking (referenced modules being corrected
        or retracted),
  \item Terminology drifting (the meaning of terms shifting),
  \item Accessibility decreasing (the module's potential extensions
        shrinking as surrounding modules evolve without it).
\end{itemize}

Semantic Yarncrawlers are agents that continuously traverse the
knowledge graph, identifying degraded connections, updating outdated
propositions, resolving terminological drift, and maintaining the
admissibility structures of modules. Large language models — if
equipped with accurate models of their own uncertainty and
well-designed update mechanisms — could serve as semantic Yarncrawlers.

\begin{namedthm}{The Yarncrawler Principle}
  Any complex system that must maintain coherence over time in a
  changing environment requires distributed repair agents that:
  (1) are always embedded in the system,
  (2) act on local information,
  (3) coordinate stigmergically through environmental modification,
  (4) optimize for accessibility entropy rather than immediate performance.
  Systems without Yarncrawlers accumulate constraint violations until
  catastrophic repair or collapse becomes necessary.
\end{namedthm}

\begin{exercise}
  Model a simple road repair network with 3 Yarncrawlers on a 10-node
  graph. Define $\kappa_I$ as the age of each road segment since last
  repair. Define a simple adaptive strategy $\pi_r$ and simulate the
  dynamics. Compare total constraint violation over time to a scheduled
  maintenance strategy that repairs each node at fixed intervals.
\end{exercise}

\begin{exercise}
  Wikipedia can be modeled as a stigmergic semantic maintenance system.
  Identify the Yarncrawler agents (human editors, bots), the environmental
  modifications they make (edits, reversions, talk page comments),
  and the stigmergic signals they leave (edit history, watchlists,
  talk page flags). Where does this model succeed and where does it
  break down?
\end{exercise}

\begin{exercise}[title={Design}]
  Design a Yarncrawler system for a university library. What constitutes
  ``degradation'' for a library's semantic infrastructure? What actions
  can a Yarncrawler take? What stigmergic signals are available?
  How does the system balance maintenance of existing collections
  against acquisition of new materials?
\end{exercise}
